\section{Parallel \cpi{} Construction}\label{sec:parabuild}

Once \spinStar{} makes peeling counter-only, \cpi{} construction becomes the load-bearing cost of the pipeline (Figure~\ref{fig:breakdown}).
%
This bottleneck shift is why construction is part of the algorithmic pipeline rather than a detached implementation detail.
%
We observe that the static-\cpi{} invariant makes construction embarrassingly parallel.
%
With no mutation contract crossing threads and each $s$-clique uniquely owned by its smallest-rank seed, the standard degeneracy-ordered Bron--Kerbosch template runs lock-free across $T$ threads.

\stitle{Why the seed loop can be parallel.}
The \cpi{} builder follows the standard degeneracy-ordered Bron--Kerbosch search with Tomita pivoting~\cite{NuclearCD,BronKerbosch,eppstein,tomita}.
%
Each $s$-clique has one owner: the smallest vertex of the clique in the degeneracy order.
%
Therefore, the outer loop over seed vertices partitions the emitted $s$-clique families.
%
No two seeds own the same $s$-clique.

Static CPI also changes what construction must produce.
%
The mutable baseline builds residual tree state because later peeling may edit that state.
%
Our $r=1$ pipeline never edits the stored hold or pivot sets.
%
The builder therefore only has to emit the path set $\mathcal{P}$; no per-path mutation contract has to be maintained across the parallel boundary.

\stitle{Independent per-thread enumeration.}
\paraBuild{} takes $G$, $s$, and a thread count $T$.
%
Each thread owns a contiguous range of seed vertices and runs the same seed-rooted enumeration as the sequential builder on each of its seeds, accumulating the emitted paths into its own output set $\mathcal{P}_t$.
%
The degeneracy ordering gives every $s$-clique a unique smallest-rank owner, so the sets $\{\mathcal{P}_t\}_t$ are pairwise disjoint; the parallel output is therefore the disjoint union $\mathcal{P}=\bigcup_t\mathcal{P}_t$, and the merge is purely the set union, with no synchronisation across threads during enumeration and no deduplication or rewrite after.
%
The expensive work is the Bron--Kerbosch search under each seed; the merge contributes no per-path arithmetic.

\begin{algorithm}[t]
\algfontsize
\caption{\paraBuild{}: parallel \cpi{} construction.}
\label{alg:parabuild}
\KwIn{Graph $G=(V,E)$; clique size $s$; thread count $T$.}
\KwOut{Static \cpi{} $\mathcal{P}$.}

Compute a degeneracy ordering $(v_1,\ldots,v_n)$ of $G$ and orient every edge from lower to higher rank, giving $N^{+}(\cdot)$\;
Partition the seed sequence $(v_1,\ldots,v_n)$ into $T$ contiguous chunks $C_1,\ldots,C_T$\;
\ForEach{thread $t\in\{1,\ldots,T\}$ \emph{in parallel}}{
    $\mathcal{P}_t \gets \emptyset$\;
    \ForEach{seed vertex $v \in C_t$}{
        $\textsc{ExpandPath}_t(N^{+}(v),\ \{v\},\ \emptyset)$ \tcp*{one root per vertex}
    }
}
$\mathcal{P} \gets \bigcup_{t=1}^{T} \mathcal{P}_t$\;
\Return{$\mathcal{P}$}
\SetKwFunction{Expand}{ExpandPath\textsubscript{$t$}}
\SetKwProg{Fn}{Function}{:}{}
\Fn{\Expand{$P,\ \Vh,\ \Vp$}}{
    \lIf{$|\Vh| > s$}{\Return}
    \If{$P = \emptyset$}{
        \lIf{$|\Vh| + |\Vp| \ge s$}{
            append the path $(\Vh,\Vp)$ to $\mathcal{P}_t$
        }
        \Return
    }
    $u_{\mathrm{piv}} \gets \arg\max_{x\in P}\, |P \cap N^{+}(x)|$ \tcp*{Tomita pivot}
    \ForEach{$v \in P \setminus N(u_{\mathrm{piv}})$}{
        $P \gets P \setminus \{v\}$\;
        \uIf{$v = u_{\mathrm{piv}}$}{
            \Expand{$P\cap N(v),\ \Vh,\ \Vp\cup\{v\}$}\;
        }
        \Else{
            \Expand{$P\cap N(v),\ \Vh\cup\{v\},\ \Vp$}\;
        }
    }
}
\end{algorithm}

Algorithm~\ref{alg:parabuild} mirrors Algorithm~\ref{alg:buildcpi} line-for-line: Line~1 computes the same degeneracy ordering and edge orientation, and the function body $\textsc{ExpandPath}_t$ at the bottom is identical to Algorithm~\ref{alg:buildcpi}'s $\textsc{ExpandPath}$ except that it appends to the thread-local output $\mathcal{P}_t$ instead of the global $\mathcal{P}$.
%
The only structural change is at the driver loop: Algorithm~\ref{alg:buildcpi}'s sequential \emph{for} over seeds is replaced by Line~2 partitioning the seed sequence into $T$ chunks and Lines~3--7 running one chunk per thread, with Line~8 merging the thread-local outputs.
%
Every $s$-clique has a unique smallest-rank owner under the degeneracy ordering, so the chunks $C_1,\ldots,C_T$ induce a partition of $\mathcal{P}$: each path is emitted by exactly one thread, so the per-thread outputs are pairwise disjoint and $\bigcup_t \mathcal{P}_t$ is the same set of paths Algorithm~\ref{alg:buildcpi} would have produced sequentially.
%
The choice of $T$, the partition of seeds, and the in-memory layout used to store $\mathcal{P}$ are implementation knobs; they affect throughput but not the output set.

\stitle{Correctness.}
\paraBuild{} is a parallel schedule of the same seed-owned enumeration.
%
Seed ownership gives disjoint emitted $s$-clique families.
%
The deterministic merge changes only the storage order, not the represented path families.
%
The resulting path set satisfies Definition~\ref{def:path}.
%
Therefore, the support formula of Lemma~\ref{lem:supp} applies unchanged.

\stitle{Scope.}
The design parallelizes construction because construction is the phase exposed by the static-\cpi{} pipeline.
%
It does not claim a lower bound showing that \spinStar{} cannot be parallelized.
%
It also does not claim ideal scaling.
%
The scaling evidence is reported in Figure~\ref{fig:par}.
